\documentclass[main.tex]{subfiles}

\begin{document}
\section{TikZ}

\begin{frame}[fragile]{TikZ (ist kein Zeichenprogramm)}
    \begin{itemize}
        \item Mit TikZ wird nicht gezeichnet, es wird programmiert
        \pause
        \begin{verbatim}
    \begin{tikzpicture}
        \node at (0, 0) {Hello, World!};
    \end{tikzpicture}
        \end{verbatim}
        \vspace{-12pt}
        \pause
        \item Sehr gutes \href{https://ctan.mc1.root.project-creative.net/graphics/pgf/base/doc/pgfmanual.pdf}{Manual}: 1300 Seiten, 24 Seiten Inhaltsverzeichnis
        \medskip
        \pause
        \item Erweiterbarkeit: \verb|\usetikzlibrary{...}|
        \begin{itemize}
            \item positioning
            \item calc
            \item shapes
        \end{itemize}
    \end{itemize}
\end{frame}

\begin{frame}[fragile]{TikZ (ist kein Zeichenprogramm)}
    \begin{itemize}
        \item Befehle für das Zeichnen
        \pause
        \begin{itemize}
            \item \verb|\node|
            \item \verb|\draw|
            \item \verb|\fill|
        \end{itemize}
        \pause
        \medskip
        \item Koordinaten
        \begin{itemize}
            \item Normal: \texttt{(x, y)}
            \item Polar: \texttt{(a:b:r)}
        \end{itemize}
        \pause
        \medskip
        \item Positionierung neben anderen Nodes
        \begin{verbatim}
    \node (hello) at (0, 0) {Hello, World!};
    \node[left=of hello] {Bonjour};
        \end{verbatim}
    \end{itemize}
\end{frame}



\begin{frame}[fragile]{TikZ Optionen}
    \begin{itemize}
        \item Dem \verb|tikzpicture| können optionen mitgegeben werden
        \pause
        \medskip
        \item Im Nachhinein kann die Größe des Bildes angepasst werden
        \begin{itemize}
            \item[$\to$] \verb|\begin{tikzpicture}[scale=2]|
        \end{itemize}
        \pause
        \medskip
        \item Auch nodes sowie die meisten Befehle unterstützen das
        \begin{itemize}
            \item Rotation: \texttt{[rotate=45]}
            \item Farben: \texttt{[red]}, \texttt{[blue]}, \texttt{[orange]}
            \item Shift: \texttt{[xshift=1cm]} / \texttt{[yshift-0.5mm]}
        \end{itemize}
        \pause
        \medskip
        \item Klassischerweise: $y$-Achse im Koordinatensystem nach oben
        \item Kann geändert werden!
        \begin{itemize}
            \item[$\to$] \verb|\begin{tikzpicture}[yscale=-1]|
        \end{itemize}
        \pause
        \medskip
        \item TikZ unterstützt auch animationen!
        \begin{verbatim}
    \uncover<1-2>{\node at (0, 0) {Hello, World!};}
        \end{verbatim}
    \end{itemize}
\end{frame}


\begin{frame}[fragile]{TikZ \textit{in action}}
    \begin{onlyenv}%
        <1, 3, 5, 7, 9>
        \begin{minted}[style=autumn]{latex}
\node (biblatex) {BibLaTeX};
\node[below=6cm of biblatex] (natbib) {natbib};
        \end{minted}
    \end{onlyenv}

    \begin{onlyenv}%
        <3, 5, 7, 9>
        \begin{minted}[style=autumn]{latex}
\node[below=of biblatex] (s biblatex) {Setup};
\node[above=of natbib] (s natbib) {Setup};
        \end{minted}
    \end{onlyenv}

    \begin{onlyenv}%
        <5, 7, 9>
        \begin{minted}[style=autumn]{latex}
\node[right=of s biblatex] (1c biblatex) {1.Compile};
\node[right=of s natbib] (1c natbib) {1.Compile};
        \end{minted}
    \end{onlyenv}

    \begin{onlyenv}%
        <7, 9>
        \begin{minted}[style=autumn]{latex}
\node[right=of 1c biblatex] (biber) {Biber};
\node[right=of 1c natbib] (bibtex) {BibTeX};
        \end{minted}
    \end{onlyenv}

    \begin{onlyenv}%
    <9>
        \begin{minted}[style=autumn]{latex}
\node[right=of $(biber)!0.5!(bibtex)$] (2 compile)     {2. Compile};
        \end{minted}
    \end{onlyenv}

    \begin{onlyenv}%
        <2, 4, 6, 8, 10>
        \begin{center}
            \advance\leftskip-25pt
            \vspace{-18pt}
            \def\colorSetup{ForestGreen}
            \def\colorCompile{VioletRed}
            \def\colorBackend{DeepSkyBlue}

            \begin{tikzpicture}
                [
                block/.style={rectangle, draw, minimum height=1cm, minimum width=2cm, align=center, thick},
                init/.style={initial by arrow, initial text=},
                arrow/.style={->,>=stealth},
                every initial by arrow/.style=arrow,
                every accepting by arrow/.style=arrow,
                ]

                % Frontend Nodes
                \uncover<2->{\node[block, init] (biblatex) {BibLaTeX};}
                \uncover<2->{\node[block, init, below=6cm of biblatex] (natbib) {natbib};}

                % Setup Nodes and their descriptions
                \uncover<4->{\node[block, below=2.25cm of biblatex.west, anchor=west, draw=\colorSetup] (setup biblatex) {Setup};}

                \uncover<4->{\node[block, above=2.25cm of natbib.west, anchor=west, draw=\colorSetup] (setup natbib) {Setup};}

                % Compilation Nodes
                \uncover<6->{\node[block, right=0.7cm of setup biblatex, draw=\colorCompile] (1 compile biblatex) {1. Compile};}
                \uncover<6->{\node[block, right=0.7cm of setup natbib, draw=\colorCompile] (1 compile natbib) {1. Compile};}

                % Backend Nodes
                \uncover<8->{\node[block, right=1.25cm of 1 compile biblatex, draw=\colorBackend] (biber) {Biber};}
                \uncover<8->{\node[block, right=1.25cm of 1 compile natbib, draw=\colorBackend] (bibtex) {BibTeX};}

                % Last Compilation nodes
                \uncover<10->{\node[block, right=2.1cm of $(biber)!0.5!(bibtex)$, accepting by arrow, draw=\colorCompile] (2 compile) {2. Compile};}

            \end{tikzpicture}
        \end{center}
    \end{onlyenv}

\end{frame}

\begin{frame}[fragile]{TikZ \textit{in action}}
    \begin{onlyenv}%
        <1, 3, 5, 7>
        \begin{minted}[style=autumn]{latex}
\draw[->] (biblatex) -- (setup biblatex);
\draw (natbib) -> (setup natbib);
        \end{minted}
    \end{onlyenv}

    \begin{onlyenv}%
        <3, 5, 7>
        \begin{minted}[style=autumn]{latex}
\draw[->] (setup biblatex) --  (1c biblatex);
\draw[->] (setup natbib) -- (1c natbib);
        \end{minted}
    \end{onlyenv}

    \begin{onlyenv}%
        <5, 7>
        \begin{minted}[style=autumn]{latex}
\draw[->] (1c natbib) -- node[above] {.aux} (bibtex);
\draw[->] (1c biblatex) -- node[above] {.aux} node[below] {.bcf} (biber);
\draw[->] (1c biblatex) -- node[above right=-0.1] {.aux} (bibtex);
        \end{minted}
    \end{onlyenv}

    \begin{onlyenv}%
    <7>
        \begin{minted}[style=autumn]{latex}
\draw[->] ($(biber.east)!0.5!(biber.south east)$) -- ($(2 compile.west)!0.5!(2 compile.north west)$);
\draw[->] ($(bibtex.east)!0.5!(bibtex.north east)$) -- ($(2 compile.west)!0.5!(2 compile.south west)$);
        \end{minted}
    \end{onlyenv}

    \begin{onlyenv}%
        <2, 4, 6, 8>
        \begin{center}
            \advance\leftskip-25pt
            \vspace{-18pt}
            \def\colorSetup{ForestGreen}
            \def\colorCompile{VioletRed}
            \def\colorBackend{DeepSkyBlue}

            \begin{tikzpicture}
                [
                block/.style={rectangle, draw, minimum height=1cm, minimum width=2cm, align=center, thick},
                init/.style={initial by arrow, initial text=},
                arrow/.style={->,>=stealth},
                every initial by arrow/.style=arrow,
                every accepting by arrow/.style=arrow,
                ]

                % Frontend Nodes
                \node[block, init] (biblatex) {BibLaTeX};
                \node[block, init, below=6cm of biblatex] (natbib) {natbib};

                % Setup Nodes and their descriptions
                \node[block, below=2.25cm of biblatex.west, anchor=west, draw=\colorSetup] (setup biblatex) {Setup};

                \node[block, above=2.25cm of natbib.west, anchor=west, draw=\colorSetup] (setup natbib) {Setup};

                % Compilation Nodes
                \node[block, right=0.7cm of setup biblatex, draw=\colorCompile] (1 compile biblatex) {1. Compile};
                \node[block, right=0.7cm of setup natbib, draw=\colorCompile] (1 compile natbib) {1. Compile};

                % Backend Nodes
                \node[block, right=1.25cm of 1 compile biblatex, draw=\colorBackend] (biber) {Biber};
                \node[block, right=1.25cm of 1 compile natbib, draw=\colorBackend] (bibtex) {BibTeX};

                % Last Compilation nodes
                \node[block, right=2.1cm of $(biber)!0.5!(bibtex)$, accepting by arrow, draw=\colorCompile] (2 compile) {2. Compile};


                % Arrows
                \uncover<2->{\draw[arrow] (biblatex) -- (setup biblatex);}
                \uncover<2->{\draw[arrow] (natbib) -- (setup natbib);}

                \uncover<4->{\draw[arrow] (setup biblatex) --  (1 compile biblatex);}
                \uncover<4->{\draw[arrow] (setup natbib) -- (1 compile natbib);}

                \uncover<6->{\draw[arrow] (1 compile biblatex) -- node[above] {.aux} node[below] {.bcf} (biber);}
                \uncover<6->{\draw[arrow] (1 compile biblatex) -- node[above right=-0.1] {.aux} (bibtex);}
                \uncover<6->{\draw[arrow] (1 compile natbib) -- node[above] {.aux} (bibtex);}

                \uncover<8->{\draw[arrow] ($(biber.east)!0.5!(biber.south east)$) -- node[above right=0.1cm, xshift=-0.3cm] {.bbl} ($(2 compile.west)!0.5!(2 compile.north west)$);}
                \uncover<8->{\draw[arrow] ($(bibtex.east)!0.5!(bibtex.north east)$) -- node[below right=0.1cm, xshift=-0.3cm] {.bbl} ($(2 compile.west)!0.5!(2 compile.south west)$);}

            \end{tikzpicture}
        \end{center}
    \end{onlyenv}

\end{frame}



\def\makeNode#1#2{

    \draw[rounded corners] (#2, 0) rectangle (#2+\nodeWidth, \nodeHeight);
    \path (#2, \nodeHeight) -- (#2+\nodeWidth, \nodeHeight) node[midway, below=0.27] (Text #1) {Node #1};
    \node[ellipse, draw, below=0.5cm of Text #1] (Prozess #1) {Prozess};

    \begin{scope}[every node/.append style={rounded corners}]
        \node[rectangle, draw, below=0.5cm of Prozess #1] (IPC #1) {IPC API};
        \node[rectangle, draw, below=1cm of IPC #1, xshift=-13*#2, minimum width=2.5cm, fill=white] (NIC #1) {NIC};
        \node[rectangle, draw, above=-0.1cm of NIC #1, xshift=10*#2, minimum width=1cm, fill=white] (Queue #1) {Queue};
    \end{scope}

    \node[below=2.2cm of Prozess #1, xshift=5*#2] {\makecell{Operating\\ System}};
}


\begin{frame}{Was TikZ so alles kann: Intersections}
    \begin{center}
        \begin{tikzpicture}
            \def\aclr{orange}
            \def\bclr{red}

            \draw[DeepSkyBlue] (0, 0) circle (2.5);
            \draw[rotate=-45, draw=\aclr] (-1, 0) ellipse (1.2 and 1.7);
            \draw[rotate=45, draw=\bclr] (-0.5, -0.3) ellipse (1.5 and 1.9);

            \begin{scope}
                \clip[rotate=45] (-0.5, -0.3) ellipse (1.5 and 1.9);
                \clip[rotate=-45] (-1, 0) ellipse (1.2 and 1.7);
                \draw[rotate=-45, fill=\aclr!50!\bclr, opacity=0.7] (-1, 0) ellipse (1.2 and 1.7);
            \end{scope}

            \node[\aclr] (A) at (-3, 1.7) {$A$};
            \draw[bend left=20, -Latex, \aclr] (A) to (-0.5, 1.7);

            \node[\bclr] (B) at (3, -1.7) {$B$};
            \draw[bend right=20, -Latex, \bclr] (B) to (0.5, -1.7);

        \end{tikzpicture}
    \end{center}
\end{frame}


\begin{frame}{Was TikZ so alles kann: Prozessdiagramme}
    \begin{center}
        \advance\leftskip-7pt
        \def\nodeWidth{3.5}
        \def\nodeHeight{5.5}
        \def\theBuffer{3.5}

        \begin{tikzpicture}
            [every path/.append style={->, shorten >= 1pt}]

            \makeNode{1}{-\theBuffer}
            \makeNode{2}{\theBuffer}

            \draw[out=185, in=175, looseness=2.325] (Prozess 1) to node[left, midway] {} (IPC 1);
            \draw[out=185, in=175, looseness=2.5] (IPC 1) to (Queue 1);

            \draw (NIC 1) to node[midway, above=0.5] (m) {\textcolor{ForestGreen}{Information\mail}} (NIC 2);

            \draw[out=5, in=-5, looseness=2.325] (IPC 2) to node[right, midway] {}(Prozess 2);
            \draw[out=5, in=-5, looseness=2.5] (Queue 2) to(IPC 2);
        \end{tikzpicture}

        \vspace{10pt}
        Bonuspunkte, wenn ihr wisst wo dieses Bild herkommt ;)
    \end{center}
\end{frame}

\begin{frame}{Was TikZ so alles kann: Fraktale}
    \begin{center}
        \begin{tikzpicture}[decoration=Koch curve type 2]
            \only<1>{\draw decorate{ (-5,0) -- (5,0) };}
            \only<2>{\draw decorate{ decorate{ (-5,0) -- (5,0) }};}
            \only<3>{\draw decorate{ decorate{ decorate{ (-5,0) -- (5,0) }}};}
%            \only<4>{\draw decorate{ decorate{ decorate{ decorate{ (-5,0) -- (5,0) }}}};}
        \end{tikzpicture}
    \end{center}

\end{frame}


\begin{frame}{Was TikZ so alles kann: 3d Bilder}
    % Graciously adopted from the TikZ manual
    \begin{tikzpicture}
        [z={(10:10mm)},x={(-45:5mm)}]
        \def\wave{
            \draw[fill,thick,fill opacity=.2]
            (0,0) sin (1,1) cos (2,0) sin (3,-1) cos (4,0)
            sin (5,1) cos (6,0) sin (7,-1) cos (8,0)
            sin (9,1) cos (10,0)sin (11,-1)cos (12,0);
            \foreach \shift in {0,4,8}
                {
                \begin{scope}[xshift=\shift cm,thin]
                    \draw (.5,0)  -- (0.5,0 |- 45:1cm);
                    \draw (1,0)   -- (1,1);
                    \draw (1.5,0) -- (1.5,0 |- 45:1cm);
                    \draw (2.5,0) -- (2.5,0 |- -45:1cm);
                    \draw (3,0)   -- (3,-1);
                    \draw (3.5,0) -- (3.5,0 |- -45:1cm);
                \end{scope}
            }
        }
        \begin{scope}[canvas is zy plane at x=0,fill=blue]
            \wave
            \node at (6,-1.5) [transform shape] {magnetic field};
        \end{scope}
        \begin{scope}[canvas is zx plane at y=0,fill=red]
            \draw[help lines] (0,-2) grid (12,2);
            \wave
            \node at (6,1.5) [rotate=180,xscale=-1,transform shape] {electric field};
        \end{scope}
    \end{tikzpicture}
\end{frame}

\begin{frame}{Was TikZ so alles kann: Mindmaps}
    \begin{center}
        \advance\leftskip-20pt

        % Graciously adopted from the TikZ manual
        \begin{tikzpicture}
            \path[mindmap,concept color=black,text=white]
            node[concept] {\makecell{Computer\\Science}}
            [clockwise from=0]
            % note that the `concept color' is passed to the `child'(!)
            child[concept color=blue] {
                node[concept] {applied}
                [clockwise from=-30]
            }
            % note that `sibling angle' can only be defined in
            % `level 1 concept/.append style={}'
            child[concept color=green!50!black] {
                node[concept] {practical}
                [clockwise from=35]
                child { node[concept] {algorithms} }
                child { node[concept] {data structures} }
                child { node[concept] {pro\-gramming languages} }
                child { node[concept] {software engineer\-ing} }
            }
            child[concept color=red] { node[concept] {technical} }
            child[concept color=orange] { node[concept] {theoretical} };
        \end{tikzpicture}

    \end{center}
\end{frame}


\end{document}
